\section{Related Work}

\noindent\textbf{AI Agents.} The rapid advancement of LLMs has sparked significant interest in autonomous AI agents capable of planning and using tools. Early general-purpose agents such as ReAct \citep{yao2022react}, HuggingGPT \citep{shen2023hugginggpt}, and AutoGen \citep{wu2024autogen} leverage external tools to break down and execute complex, open-ended tasks. In the field of software development, systems such as Voyager \citep{wang2023voyager}, AlphaCode \citep{li2022competition}, SWE-agent \citep{yang2024swe}, and Claude Code \citep{liu2026dive} use interfaces that recognize execution feedback and the environment to iteratively debug and solve coding problems. Recently, this paradigm has expanded into the fields of automated machine learning (ML) engineering and data science. Frameworks such as MLAgentBench \citep{huang2023mlagentbench}, OpenHands \citep{wang2025openhands}, AIDE \citep{jiang2025aide}, MLE-STAR \citep{nam2026mle}, and MARS \citep{chen2026mars} automate end-to-end ML workflows by executing and improving modeling pipelines, while data science agents such as DA-Agent \citep{huang2024code}, Data Interpreter \citep{hong2025data}, and DS-STAR \citep{nam2026ds} tackle heterogeneous data challenges through structured planning and recursive verification. In this paper, we introduce \sname, a specialized AI agent for autonomous research, from idea generation and implementation to the writing of high-quality publishable academic papers.

\noindent\textbf{Autonomous Research Agents.} Autonomous research agents have rapidly evolved from ML problems into sophisticated, multi-stage pipelines that coordinate the entire scientific workflow—from literature-based analysis and hypothesis generation to execution and drafting of research papers. Early end-to-end pipelines, such as AI Scientist \citep{lu2024ai}, established this automation paradigm but suffered from execution instability and issues with hallucinated writing; subsequent versions, such as AI Scientist-v2 \citep{yamada2025ai}, mitigated these problems through a best-first-tree search for experimental branches and review-based reporting. At the same time, modular systems have focused on resolving specific bottlenecks; for example, PaperOrchestra \citep{song2026paperorchestra} compiles unconstrained experiment logs and ideas into LaTeX manuscripts, while ScholarPeer \citep{goyal2026scholarpeer} deploys a multi-agent system that acts as peer-reviewers. Other frameworks introduce human-in-the-loop gates \citep{schmidgall2025agent}, apply evolutionary optimization to algorithmic discovery \citep{novikov2025alphaevolve, lyu2026evoscientist}, or focus on single-metric optimization for a target codebase \citep{li2026autosota, liu2026autoresearchclaw, tang2026ai, weng2025deepscientist}; however, they are generally limited to optimizing a single scalar metric on a single dataset. Crucially, even verifiability-centric systems like ScientistOne \citep{meng2026scientistone} lack the ability to perform an active analytical expansion. Our \sname~addresses these fundamental gaps by conducting systematic experiments on various benchmark datasets and incorporating a dynamic peer-review loop, automating the iterative process of performing complementary ablation experiments, refining ideas, and revising manuscripts, thus producing expert-level papers that meet the rigorous standards of top academic venues.